\documentclass[12pt, a4paper]{article}

% Pacchetti
\usepackage[utf8]{inputenc}
\usepackage[T1]{fontenc}
\usepackage[italian]{babel}
\usepackage{geometry}
\usepackage{titlesec}
\usepackage{enumitem}
\usepackage{fancyhdr}
\usepackage{xcolor}
\usepackage{booktabs}
\usepackage{array}
\usepackage{parskip}
\usepackage{hyperref}
\usepackage{amssymb}

% Margini
\geometry{top=2.5cm, bottom=2.5cm, left=3cm, right=3cm}

% Colori
\definecolor{primary}{RGB}{30, 60, 120}
\definecolor{lightgray}{RGB}{245, 245, 245}
\definecolor{darkgray}{RGB}{80, 80, 80}
\definecolor{successgreen}{RGB}{30, 120, 60}
\definecolor{alertred}{RGB}{180, 30, 30}
\definecolor{warningyellow}{RGB}{150, 100, 0}

% Stile titoli sezioni
\titleformat{\section}
  {\large\bfseries\color{primary}}
  {\thesection.}{0.5em}{}[\titlerule]

\titleformat{\subsection}
  {\normalsize\bfseries\color{darkgray}}
  {\thesubsection.}{0.5em}{}

% Header e Footer
\pagestyle{fancy}
\fancyhf{}
\rhead{\small\textcolor{darkgray}{Analisi della Trascrizione}}
\lhead{\small\textcolor{darkgray}{MirrorFish Trading Offline}}
\rfoot{\small\textcolor{darkgray}{\thepage}}
\renewcommand{\headrulewidth}{0.4pt}

% Hyperlink
\hypersetup{
  colorlinks=true,
  linkcolor=primary,
  urlcolor=primary
}

% ─────────────────────────────────────────────
\begin{document}

% Intestazione / Titolo
\begin{center}
  {\LARGE\bfseries\color{primary} ANALISI DELLA TRASCRIZIONE}\\[0.4cm]
  {\large Progetto: \textit{MirrorFish Trading Offline}}\\[0.2cm]
  \textcolor{darkgray}{\rule{0.6\textwidth}{0.4pt}}
\end{center}

\vspace{0.8cm}

% ─────────────────────────────────────────────
\section{Analisi Preliminare}

La trascrizione è estremamente frammentata, incoerente e utilizza un linguaggio che sembra mescolare terminologia tecnica (programmazione, strutture dati) con termini altamente criptici o incomprensibili (es.\ \textit{``d'attabeditoriale,''} \textit{``latine,''} \textit{``graff''}). Il contesto generale è difficile da determinare, ma l'argomento centrale sembra ruotare attorno a un processo di \textbf{trasformazione di dati o di funzioni}.

\vspace{0.5cm}

% ─────────────────────────────────────────────
\section{Riassunto Generale}

La conversazione, sebbene difficile da decifrare a causa della formulazione sintattica e semantica erratica, riguarda l'introduzione o la spiegazione di una \textbf{funzione specifica} che opera attraverso una trasformazione di dati. Si menziona un processo chiamato \textit{``trasformazione dei conti''} che risulta in un formato o struttura dati chiamato \textbf{``Graff''} (descritto come un \textit{``dattabeditoriale a forma di graffo per calità''}). L'obiettivo finale sembra essere mostrare come scrivere questa funzione in pratica (\textit{``Inveito questo è il funzione che bisogna scrivere per dattere''}).

\vspace{0.3cm}

\noindent\textcolor{warningyellow}{\textbf{$\triangle$ Nota di Incertezza:}} Gran parte del linguaggio è molto probabilmente basato su termini mal pronunciati o altamente specialistici/erronei. È difficile stabilire se si tratti di concetti reali di programmazione o se ci sia un'interruzione tra diversi argomenti.

\vspace{0.5cm}

% ─────────────────────────────────────────────
\section{Decisioni Prese}

Non è possibile identificare decisioni prese all'interno del frammento fornito. La conversazione è un flusso esplicativo o dimostrativo, non un momento di confronto che porta a un accordo o a una scelta definitiva.

\vspace{0.5cm}

% ─────────────────────────────────────────────
\section{Task / Azioni da Fare}

\begin{enumerate}[leftmargin=1.5cm, label=\textbf{\arabic*.}]

  \item \textbf{Riprovare l'elaborazione del contenuto}\\
  È necessario chiedere al relatore di riformulare la spiegazione utilizzando un linguaggio più semplice e focalizzato sul concetto tecnico che vuole trasmettere (es.\ sostituire \textit{``dattabeditoriale''} e \textit{``Graff''} con termini standard).

  \item \textbf{Fornire esempi concreti}\\
  Se l'argomento è tecnico, è essenziale che vengano forniti snippet di codice o esempi di input/output per rendere chiara la trasformazione descritta.

\end{enumerate}

\vspace{0.5cm}

% ─────────────────────────────────────────────
\section{Punti Importanti o Critici}

\subsection{\textcolor{successgreen}{Punti di Forza — Concetti Riconoscibili}}

\begin{itemize}[leftmargin=1.5cm]
  \item \textbf{Processo di Trasformazione:} Il concetto fondamentale è che A viene trasformato in B (trasformazione dei conti $\rightarrow$ dattabeditoriale chiamato Graff).
  \item \textbf{Obiettivo Codificabile:} Vi è l'intenzione di scrivere una ``funzione'' specifica.
\end{itemize}

\subsection{\textcolor{alertred}{Punti Critici e Ambiguità — Da Verificare Urgentemente}}

\begin{itemize}[leftmargin=1.5cm]
  \item \textbf{Terminologia Specifica:} I termini \textit{``d'attabeditoriale,''} \textit{``latine,''} \textit{``graff''} e \textit{``dattabeditoriale''} non sono standard nel linguaggio tecnico o generale. È cruciale identificare il significato corretto di questi nomi.
  \item \textbf{Confusione Concettuale:} La transizione tra \textit{``vi dello uso,''} \textit{``i fai della fondamentale battalandosi''} e la discussione del \textit{``Graff''} è troppo brusca, suggerendo un possibile \textbf{disallineamento del soggetto} da parte del relatore.
  \item \textbf{Flusso Logico Interrotto:} Il report è estremamente difficile da seguire perché il relatore sembra saltare tra descrizioni astratte (trasformazione) e specifici nomi tecnici (Graff) senza costruire un ponte logico chiaro tra i passaggi.
\end{itemize}

\vspace{0.8cm}
\noindent\rule{\textwidth}{0.4pt}

\vspace{0.3cm}
\noindent\textbf{Sintesi dell'Analista:} La trascrizione richiede un intervento esterno per la pulizia del linguaggio. Il focus deve essere spostato dai termini criptici alla funzione matematica o logica che il relatore vuole dimostrare.

\end{document}